\section{Related Work}
\label{sec:related}
\vspace{-1mm}

\paragraph{EEG Foundation Models.}
The development of large-scale pretrained models for EEG has accelerated rapidly. LaBraM~\cite{jiang2024labram} introduced the first large-scale EEG foundation model, pretraining on 2,500 hours of data from 20 datasets using vector-quantized neural spectrum prediction. CBraMod~\cite{wang2025cbramod} proposed a criss-cross transformer that factorizes spatial and temporal attention, pretrained on the Temple University Hospital EEG Corpus (TUEG). REVE~\cite{elouahidi2025reve} scaled pretraining to 60,000 hours across 25,000 subjects with a novel 4D positional encoding scheme. CoMET~\cite{li2025comet} combined contrastive learning with masked autoencoding, explicitly separating mask tokens from the encoder (unlike CBraMod's grid-constrained approach). EEG-DINO~\cite{wang2025eegdino} applied hierarchical self-distillation to learn EEG representations without explicit reconstruction targets. mdJPT~\cite{mdjpt2025} explored multi-dataset joint pretraining for cross-domain EEG generalization. On the adaptation side, EEG-GraphAdapter~\cite{eeg_graphadapter2024} proposed adapter-based graph learning for parameter-efficient fine-tuning of EEG foundation models. These models achieve state-of-the-art performance across diverse EEG tasks but treat emotion recognition as a standard classification problem with integer labels. In contrast, \ours leverages the relational structure among emotions as an additional supervisory signal, complementing rather than replacing the pretrained representations.

\paragraph{Emotion Representations in LLMs.}
Recent interpretability research has revealed that LLMs develop structured internal representations of emotion concepts. Sofroniew et al.~\cite{sofroniew2026anthropic} demonstrated that Claude's emotion vectors organize along valence-arousal axes, are causally functional (driving behavior such as reward hacking), and persist across diverse contexts. Tak et al.~\cite{tak2025emotion_interpretability} showed through probing and activation patching that emotion processing concentrates in mid-layer multi-head self-attention, with peak encoding at 50--75\% depth. Layer-wise probing of LLMs for emotion has revealed depth-dependent encoding of affective information~\cite{decoding_emotion_deep2024}, while controllable generation work has demonstrated that emotion vectors can steer text output~\cite{emotion_vectors_controllable2025}. Brain--LLM alignment studies find that scaling, but not instruction tuning beyond a threshold, increases representational correspondence with neural data~\cite{scaling_brain_alignment2024}. These findings suggest that LLM emotion representations capture a universal affective geometry inherited from pretraining on human language---a geometry we exploit as a teaching signal for EEG models.

\paragraph{EEG Emotion Recognition with Language Models.}
Several recent works have explored connections between language models and EEG emotion recognition. EmotionCLIP~\cite{emotionclip2025} reformulates emotion recognition as EEG-text matching using frozen CLIP text encoders with prompt templates, achieving strong cross-subject performance on SEED. EMOD~\cite{emod2026} projects discrete emotion labels into a continuous valence-arousal (V-A) space and applies soft-weighted supervised contrastive learning, pretraining on 8 EEG datasets. E$^2$-LLM~\cite{ma2026e2llm} directly integrates an EEG encoder with Qwen-based LLMs through multi-stage training including chain-of-thought reasoning. Our approach differs fundamentally: rather than using LLM text embeddings of emotion \emph{labels} (EmotionCLIP) or predefined V-A coordinates (EMOD), we extract emotion vectors from LLM \emph{hidden states} during emotion-laden story generation, capturing richer relational structure that emerges from the model's internal representations.

\paragraph{Knowledge Distillation for EEG.}
Knowledge distillation~\cite{hinton2015distilling} has been applied to EEG primarily for model compression---transferring knowledge from high-density to low-density electrode configurations. Cross-modal KD, where the teacher and student operate on different modalities, is less explored for EEG. CMCRD~\cite{cmcrd2025} proposed cross-modal contrastive representation distillation for bridging heterogeneous feature spaces. In our setting, the LLM teacher provides a \emph{relational} supervision signal (similarity structure among emotion classes) rather than transferring a specific representation, making our approach closer to relational knowledge distillation than standard feature-matching KD.

\paragraph{Parameter-Efficient Fine-Tuning for Neural Signals.}
PEFT methods including adapters~\cite{houlsby2019adapters} and LoRA~\cite{hu2022lora} have transformed NLP and vision fine-tuning but remain underexplored for EEG foundation models. We systematically compare bottleneck adapters, LoRA, and FiLM conditioning~\cite{perez2018film} for EEG emotion recognition, finding that adapters with frozen backbone not only provide parameter efficiency but also act as a regularizer that prevents degradation of pretrained representations---a finding with implications for the broader EEG foundation model community.
